% ===========================================================================
\section{Affine Varieties}\label{sec:affine}
\warmth{0}\quad\whisper{The geometry side, and the bridge. ``Geometry is but drawn algebra.''}

\begin{strand}
A \keyword{variety} is the solution set of polynomial equations,
$\V(I)=\{p\in K^n: f(p)=0\ \forall f\in I\}$. Different ideals can give the same variety, so geometry
sees only $I$ \emph{up to radical}. The reverse map $\Ideal(W)=\{f:f|_W=0\}$ sends a shape back to its
(radical) ideal. These two maps are the dictionary; this section makes them precise.
\end{strand}

\block{The two maps $\V$ and $\Ideal$}

\begin{definition}[variety and vanishing ideal]\label{def:VandI}
For an ideal $I\subseteq K[\xb]$ and a subset $W\subseteq K^n$,
\[
  \V(I)=\{p\in K^n: f(p)=0\ \forall f\in I\},\qquad
  \Ideal(W)=\{f\in K[\xb]: f(p)=0\ \forall p\in W\}.
\]
$\Ideal(W)$ is always a \keyword{radical} ideal. A subset is a \emph{variety} iff $W=\V(\Ideal(W))$.
\end{definition}

\recall{Both maps reverse inclusions. Fill: $I\subseteq J\Rightarrow \V(I)\,\fillin[1cm]\,\V(J)$, and $V\subseteq W\Rightarrow\Ideal(W)\,\fillin[1cm]\,\Ideal(V)$.}

\block{Irreducible $=$ prime: the keystone}

\begin{definition}[irreducible]
A variety is \keyword{irreducible} if it is not a finite union of proper subvarieties:
$\V(I)=\V(J)\cup\V(J')\Rightarrow\V(I)=\V(J)$ or $\V(J')$.
\end{definition}

\begin{proposition}[Prop.\ 2.3]\label{prop:irr-prime}
A variety $W\subseteq K^n$ is irreducible \;$\iff$\; its ideal $\Ideal(W)$ is \fillin[2.5cm].
\end{proposition}
\begin{proof}
($\Leftarrow$) Suppose $\Ideal(W)$ prime and $W=\V(J)\cup\V(J')$ with $W\ne\V(J)$. Pick $f\in J$,
$v\in W$ with $f(v)\ne0$, so $f\notin\Ideal(W)$. For $g\in J'$, the product $fg$ vanishes on
$\V(J)\cup\V(J')=W$, so $fg\in\Ideal(W)$; primality gives $g\in\Ideal(W)$. Hence $W\subseteq\V(J')$.
\TODO{($\Rightarrow$) $W$ irreducible: if $fg\in\Ideal(W)$ then $W=(W\cap\V(f))\cup(W\cap\V(g))$; conclude $f$ or $g\in\Ideal(W)$}
\end{proof}

\begin{strand}
This is the dictionary's keystone: \keyword{prime ideals} $\leftrightarrow$ \keyword{irreducible
varieties}, exactly as prime numbers are to $\Z$. In applications (algebraic statistics), a model is
the image of a polynomial map; its Zariski closure is irreducible, so its ideal is prime, and is
computed as the kernel of the dual ring map. The closed sets of the \keyword{Zariski topology} are
precisely the varieties; the coordinate ring $K[W]=K[\xb]/\Ideal(W)$ is the ring of polynomial
functions on $W$, and a map of varieties $f:W_1\to W_2$ pulls back to a ring map
$f^\ast:K[W_2]\to K[W_1]$ (a \emph{contravariant} functor).
\end{strand}

\block{Dimension and degree, geometrically}

\begin{example}[linear subspace, Ex.\ 2.15]
For a linear subspace $V\subseteq K^n$: $\dim V$ is its linear-algebra dimension, $\Ideal(V)=
\gen{x_1,\dots,x_s}$ with $s=n-\dim V$, and $\deg V=\fillin[1cm]$. \recall{Degree $1$ $\iff$ prime ideal generated by linear forms.}
\end{example}

A method to get $\dim\V(I)$: compute $\inn_\prec(I)$; find the smallest variable set $S$, $|S|=d$, hitting
every generator; then $\dim=n-d$.

\begin{theorem}[Bézout, Thm.\ 2.16]\label{thm:bezout}
Let $f_1,\dots,f_k$ be \emph{general} polynomials of degrees $d_1,\dots,d_k>0$ in $n$ variables, and
$I=\gen{f_1,\dots,f_k}$. Then
\[
  \dim(I)=\fillin[1.5cm],\qquad \deg(I)=\fillin[2cm].
\]
(Such $I$ is a \keyword{complete intersection}; $\dim\ge n-k$ always, by Krull.)
\end{theorem}
\begin{yourturn}
Two general conics in the plane ($n=2,k=2,d_1=d_2=2$): Bézout predicts $\dim=\fillin[0.8cm]$ and
$\deg=\fillin[0.8cm]$ -- i.e.\ they meet in $\fillin[0.8cm]$ points (over $\overline K$). Sketch two
ellipses meeting in $4$ points to see it.
\end{yourturn}

\block{Smooth vs singular}

\begin{definition}[singular point]
$p\in\V(f)$ is \keyword{singular} if all partials vanish: $\partial f/\partial x_i(p)=0$ for all $i$.
The singular locus is $\V\!\big(\gen{f,\partial f/\partial x_1,\dots,\partial f/\partial x_n}\big)$; if
empty, $\V(f)$ is \keyword{smooth}. For an irreducible $Y$ of dimension $d$ cut by a prime $I$, the point
$p$ is singular iff $\rank\,\mathrm{Jac}(p)<n-d$ (codimension).
\end{definition}

\recall{The four nodes of the running cubic are exactly its singular locus -- the $\dim 0,\deg 4$ ideal from §\ref{sec:dimdeg}. Smoothness $=$ the tangent space has the right dimension.}

\loose{Bertini (book): an ideal generated by general polynomials is prime and its variety is smooth, provided $\dim>0$. So ``random'' equations cut out nice irreducible smooth varieties; degeneracy is special.}
